\documentclass[11pt]{article}

% ─── Packages ────────────────────────────────────────────────────────────────
\usepackage{amsmath,amssymb,amsthm}
\usepackage[margin=1in]{geometry}
\usepackage{hyperref}
\usepackage{graphicx}
\usepackage{booktabs}
\usepackage{array}
\usepackage{multirow}
\usepackage{longtable}
\usepackage{xcolor}
\usepackage{authblk}
\usepackage{setspace}
\usepackage{caption}
\usepackage{subcaption}
\usepackage{enumitem}
\usepackage{microtype}
\usepackage{siunitx}
\usepackage{pdflscape}

\hypersetup{
  colorlinks=true,
  linkcolor=blue,
  citecolor=blue,
  urlcolor=blue
}

% ─── Theorem environments ──────────────────────��─────────────────────────────
\theoremstyle{plain}
\newtheorem{theorem}{Theorem}[section]
\newtheorem{lemma}[theorem]{Lemma}
\newtheorem{proposition}[theorem]{Proposition}
\newtheorem{corollary}[theorem]{Corollary}
\theoremstyle{definition}
\newtheorem{definition}{Definition}[section]
\theoremstyle{remark}
\newtheorem{remark}{Remark}[section]

% ─── Title block ─────────────────────────────────────────────────────────────
\title{\textbf{Geometric Closure of the Design Space:\\
The Missing Chemical Classes, the Scope of the SCE Framework,\\
and the Programme for What Remains}}

\author[1]{Eric Robert Lawson}

\affil[1]{Independent Researcher, OrganismCore Initiative\\
\texttt{ORCID: 0009-0002-0414-6544}\\
Repository: \url{https://github.com/Eric-Robert-Lawson/attractor-oncology}}

\date{April 4, 2026\\
\small Pre-registration record:
\url{https://github.com/Eric-Robert-Lawson/attractor-oncology}\\
\small (timestamped commit history constitutes the priority
and pre-registration record)\\
\small Part of the SCE Framework preprint series.
Preprint~1 (foundation and fixed point):
same repository.\\
\small This paper is the series capstone: it maps the framework's
reach, its limits, and the complete programme of what remains.}

% ─── Document ────────────────────────────────────────────────────────────────
\begin{document}

\maketitle

% ─── Abstract ────────────────────────────────────────────────────────────────
\begin{abstract}
The SCE Framework (Preprints~1--5) establishes $\mathrm{SCE}^{*}
= 1.466$ as the fixed point of Li$^+$ first-shell coordination
entropy that simultaneously maximises room-temperature (RT) and
low-temperature (LT) Coulombic efficiency (CE) in lithium metal
batteries.  Five preprints derive four candidate formulations and
one arctic-optimised design, covering three orthogonal navigation
mechanisms and one temperature-responsive design principle.

The present paper closes the design space geometrically.  We
systematically enumerate the chemical classes that are \emph{absent}
from the confirmed dataset and the series candidates, derive
pre-registered SCE predictions for each class using the framework
equations and molecular structure arguments, and establish the
conditions under which each missing class is predicted to reach,
approach, or be unable to reach $\mathrm{SCE}^{*}$.

The missing classes identified and treated are:
(1)~nitrile-class solvents (acetonitrile, propionitrile),
(2)~sulfolane and sulfone-class solvents,
(3)~fluorinated ether sub-classes (HFE, BTFMD, HFTHP),
(4)~carbonate-class solvents in ether mixtures,
(5)~ionic liquid co-solvents,
and (6)~solid-state and quasi-solid electrolyte analogues.

For each class we derive: the expected SCE position relative to
$\mathrm{SCE}^{*}$, the mechanism by which it can or cannot
reach $\mathrm{SCE}^{*}$, the specific formulation prediction
where applicable, and the falsification conditions.

We also derive the framework's three fundamental scope limits:
(A)~the \emph{coordination geometry requirement} (classes where
Li$^+$ has no accessible coordination geometry distinction cannot
navigate by the geometric mechanisms), (B)~the \emph{reductive
stability requirement} (classes that are reduced at Li metal
potentials before SCE can be established are outside the framework's
operating domain), and (C)~the \emph{concentration window
requirement} (classes where SCE is concentration-invariant and
SCE$\neq\mathrm{SCE}^{*}$ at any accessible concentration are
geometrically unreachable by concentration tuning alone).

The paper concludes with the complete novelty stack of the series,
the full geometric structure of what has been derived and what
remains, and a summary of the programme for the next generation of
experimental and computational work.

The derivation record and all pre-registered conditions are archived
at \url{https://github.com/Eric-Robert-Lawson/attractor-oncology}.
\end{abstract}

\vspace{1em}
\noindent\textbf{Keywords:} lithium metal battery, electrolyte
design space, coordination entropy, SCE*, missing classes, nitrile,
sulfone, fluorinated ether, carbonate, ionic liquid, scope limits,
geometric closure, programme

\tableofcontents
\newpage

% ─── 1. Introduction ─────────────────────────────────────────────────────────
\section{Introduction}

\subsection{What the Series Has Established}

The SCE Framework series has, across five preprints:

\begin{enumerate}[leftmargin=2em]
  \item Derived $\mathrm{SCE}^{*} = 1.466$ as a fixed point of
        combined RT and LT performance from a 21-system dataset
        ($R^2 = 0.828$, $p = 4.5 \times 10^{-6}$).
  \item Established three orthogonal mechanisms for navigating
        to $\mathrm{SCE}^{*}$: ESP matching, anion reception, and
        within-class geometric diversity.
  \item Derived one temperature-responsive design principle: the
        arctic optimum.
  \item Specified six candidate formulations with pre-registered
        falsification conditions.
  \item Confirmed the fixed point against six published papers;
        none falsifies the framework.
  \item Mapped the molecular design space within the ether class
        (MU Molecular Search, same repository) and achieved
        geometric closure of the ether-class search.
\end{enumerate}

What the series has \emph{not} established is equally important:
the framework has been derived from and applied to a dataset
dominated by ether-class and phosphate ester solvents.  The
confirmed dataset contains no nitrile-class, sulfone-class,
or ionic liquid systems.  The fluorinated ether sub-classes
are represented only indirectly.  The carbonate class is absent.
Solid-state analogues are not addressed.

This paper closes the design space by addressing these absences
systematically.

\subsection{The Purpose of Geometric Closure}

A scientific framework is not complete until its scope is defined.
Scope definition requires two things: knowing what the framework
covers (the domain of confirmed application) and knowing what
it does not cover and why (the scope limits).

The SCE Framework's scope is defined by:
\begin{itemize}[leftmargin=2em]
  \item The coordination geometry requirement: the mechanism
        operates on Li$^+$ first-shell geometry.  If a solvent
        class does not produce geometrically distinct coordination
        modes under mixing, the mechanism does not apply.
  \item The reductive stability requirement: the mechanism
        operates at Li metal potentials ($\approx 0$~V vs.\
        Li/Li$^+$).  If a solvent class is reduced before the
        electrolyte can form a stable solvation structure, it is
        outside the operating domain.
  \item The concentration window requirement: the framework's
        dataset spans LiFSI concentrations of 0.5--2.0~M.  For
        classes where SCE is concentration-invariant and
        $\mathrm{SCE} \neq \mathrm{SCE}^{*}$ at all accessible
        concentrations, concentration tuning cannot reach the
        fixed point.
\end{itemize}

For each missing class, we establish which of these requirements
is satisfied, which is violated, and what consequences follow.

\subsection{The Molecular Search Closure}

The MU Molecular Search (Complete Space Closure Record, same
repository) established geometric closure of the
\emph{ether-class} design space: every ether and ether-adjacent
molecule accessible from the molecular search starting point
(MU) has been assigned to a region (reached, unreachable,
discarded) with respect to $\mathrm{SCE}^{*}$.

The present paper extends this closure to the \emph{non-ether}
classes.  The closure is not exhaustive at the molecular level
--- it is exhaustive at the \emph{class level}: every major
chemical class of Li metal battery solvent is assigned a status,
and the reasons for that status are derived from the framework.

\subsection{Structure of This Paper}

Sections~\ref{sec:nitrile}--\ref{sec:solid_state} treat each
missing class in turn.  Section~\ref{sec:scope_limits} derives
the three scope limits formally.  Section~\ref{sec:novelty_stack}
presents the complete novelty stack of the series.
Section~\ref{sec:programme} defines the programme for what
remains.  Section~\ref{sec:discussion} provides the series
synthesis.

% ─── 2. Missing Class 1: Nitriles ─────────────────────────────────────────────
\section{Missing Class 1: Nitrile-Class Solvents}
\label{sec:nitrile}

\subsection{Why Nitriles Are Absent from the Confirmed Dataset}

Nitrile solvents (acetonitrile, AN; propionitrile, PN;
glutaronitrile, GLN) are commonly used in Li-ion battery
electrolytes but are rarely used in Li \emph{metal} batteries
because:
\begin{itemize}[leftmargin=2em]
  \item AN is reductively unstable at Li metal potentials,
        forming polymerised SEI products that consume Li
        irreversibly.
  \item PN has marginally better reductive stability but remains
        problematic for extended Li metal cycling.
  \item High-dielectric nitriles (GLN) have high viscosity that
        limits ionic conductivity at low temperature.
\end{itemize}

The dataset absence is therefore a reductive stability issue,
not a coordination entropy issue.

\subsection{SCE Position of Nitrile-Class Solvents}

Nitrile oxygens --- strictly, nitrile nitrogens --- coordinate
Li$^+$ via the lone pair on nitrogen.  The C$\equiv$N--Li$^+$
interaction is linear: the nitrogen lone pair is directed along
the triple bond axis.  This produces a single, geometrically
rigid coordination mode with no rotational or conformational
sub-configurations.

\begin{proposition}[Nitrile SCE position]
LiFSI in a pure nitrile solvent has a lower SCE than any
ether-class solvent at equivalent concentration, because the
nitrile C$\equiv$N--Li$^+$ interaction is geometrically rigid:
$n_\mathrm{sig} \approx 1$ for SSIP configurations and
$n_\mathrm{sig} \approx 2$ including CIP at moderate concentration.
The maximum achievable SCE in a pure nitrile system is:
\begin{equation}
  \mathrm{SCE}(\mathrm{nitrile}, \mathrm{pure}) \leq \ln(2) = 0.693
\end{equation}
which is far below $\mathrm{SCE}^{*} = 1.466$.
\end{proposition}

\subsection{Can Nitrile/Ether Mixtures Reach SCE*?}

A nitrile/ether co-solvent system introduces a fundamentally
new coordination class: the C$\equiv$N--Li$^+$ linear
interaction vs.\ the O--Li$^+$ coordinations of the ether.
The two geometries are maximally distinct.  A nitrile/ether
mixture at appropriate ratios could in principle reach
$\mathrm{SCE}^{*}$ by cross-class ESP matching.

However, two problems arise:

\begin{enumerate}[leftmargin=2em]
  \item \textbf{ESP mismatch is too large:} The AN nitrogen
        ESP minimum ($\approx -1.2$~eV) is substantially
        lower (more negative) than the DME oxygen ESP minimum
        ($\approx -0.62$~eV).  The energy gap is $\approx 0.58$~eV,
        more than $20 k_BT$ at 298~K.  Li$^+$ overwhelmingly
        prefers AN coordination.  The population cannot be
        meaningfully split across both geometries without
        an AN mole fraction so low that AN is merely an additive.
  \item \textbf{Reductive stability:} At the AN mole fractions
        required for meaningful coordination, AN reductive
        decomposition at Li metal makes the system impractical.
\end{enumerate}

\begin{proposition}[Nitrile geometric inaccessibility]
Nitrile/ether mixtures cannot reach $\mathrm{SCE}^{*}$ by the
cross-class ESP-matching mechanism because the ESP mismatch
exceeds the population-splitting threshold:
\begin{equation}
  |\mathrm{ESP}_{\min}(\mathrm{AN}) - \mathrm{ESP}_{\min}(\mathrm{DME})|
  \approx 0.58~\mathrm{eV} \gg \delta_{\mathrm{thermal}}
\end{equation}
The nitrile class is geometrically inaccessible to the
ESP-matching mechanism.  It is also outside the reductive
stability requirement.  Classification: \textbf{Scope Limit A
and B both apply.}
\end{proposition}

\subsection{Pre-Registered Prediction}

For completeness, a pre-registered prediction: if AN is added
to DME at any volume fraction from 0.05 to 0.50 at LiFSI 1.0~M,
Raman spectroscopy will show Li$^+$ preferentially coordinated
by AN at all AN fractions above 0.05.  SCE will not exceed
$\ln(3) \approx 1.10$ at any ratio.  CE will not improve
beyond the pure DME baseline.  The ESP-driven preferential
coordination precludes population splitting.

% ─── 3. Missing Class 2: Sulfones ──────────────────────────────────────────────
\section{Missing Class 2: Sulfolane and Sulfone-Class Solvents}
\label{sec:sulfone}

\subsection{Why Sulfones Are Absent}

Sulfolane (tetramethylene sulfone, TMS) and other sulfone-class
solvents have high dielectric constants ($\varepsilon \approx 44$
for sulfolane) and moderate donor numbers (DN $\approx 14$ for
sulfolane).  They are used in Li-ion systems for their
oxidative stability but are absent from the ether-class Li metal
dataset for two reasons:
\begin{itemize}[leftmargin=2em]
  \item High melting points: sulfolane melts at $27\,^\circ$C,
        making it a solid at ambient temperature in many climates
        and severely limiting LT performance.
  \item High viscosity: sulfolane's viscosity (10.3~cP at
        30\,$^\circ$C) is 4--5$\times$ higher than DME at
        RT and increases dramatically on cooling.
\end{itemize}

The MU session (Query~6, Master Document~8) specifically
investigated HFTHP and BTFMD as fluorinated sulfone-adjacent
structures for LT performance.  The result: both showed Raman
evidence of very high CIP/AGG fractions at LiFSI 1.0~M, placing
their SCE well above $\mathrm{SCE}^{*}$.

\subsection{SCE Position of Sulfone-Class Solvents}

Sulfolane coordinates Li$^+$ via the sulfonyl oxygen (S=O).
The sulfonyl coordination is monodentate (one oxygen per
molecule).  The ESP minimum of the sulfonyl oxygen
($\approx -0.58$~eV) is similar to that of ether oxygens,
placing sulfolane in a comparable ESP range to DME.

\textbf{The ESP finding of the MU session:} The MU search
(Sulfolane ESP Finding --- Noted But Discarded section of the
Molecular Search Closure Record) found that sulfolane's ESP
minimum is close to DME's ESP minimum ($|\Delta\mathrm{ESP}|
\approx 0.04$~eV $\approx 1.5 k_BT$).  This places sulfolane
in the within-class geometric diversity regime relative to ether
solvents: it is not ESP-matched within $k_BT$, but it is close
enough that a sulfolane/ether mixture might produce meaningful
population splitting.

\begin{proposition}[Sulfolane SCE range]
LiFSI in sulfolane has $\mathrm{SCE} > \mathrm{SCE}^{*}$
at moderate concentrations (0.5--1.0~M) due to the high
dielectric constant promoting CIP/AGG formation and high
$n_\mathrm{sig}$.  Sulfolane sits above $\mathrm{SCE}^{*}$
at RT.  Adding an ether co-solvent can reduce SCE toward
$\mathrm{SCE}^{*}$ from above --- the anion-reception or
geometric ordering direction.
\end{proposition}

\subsection{Predicted Mechanism: Ordering from Above}

Sulfolane/DME mixtures approach $\mathrm{SCE}^{*}$ from
\emph{above}: sulfolane's high dielectric environment promotes
CIP/AGG, raising SCE above the fixed point.  Adding DME
dilutes the CIP/AGG fraction (lower dielectric constant,
bidentate coordination that competes for the Li$^+$ primary
shell) and reduces SCE toward $\mathrm{SCE}^{*}$.

This is the same direction as the anion receptor mechanism
(Preprint~3) but operated by dielectric dilution rather than
Lewis acid sequestration.

\textbf{Pre-registered prediction (sulfolane/DME):}
LiFSI 1.0~M in sulfolane/DME at 1:4 v/v will show:
\begin{itemize}[leftmargin=2em]
  \item Raman SCE $\approx 1.45$--$1.50$ (approaching
        $\mathrm{SCE}^{*}$ from above).
  \item CE$_\mathrm{RT}$ between $73\%$ and $78\%$.
  \item CE$_\mathrm{LT}$ between $58\%$ and $62\%$.
\end{itemize}

\textbf{Limiting factor:} Sulfolane's melting point ($27\,^\circ$C)
constrains the range of sulfolane/DME ratios accessible at LT.
At sulfolane fractions above $\approx 0.25$ v/v in DME, the
mixture may partially solidify at $-20\,^\circ$C.  The
prediction applies to the 1:4 ratio only.

\textbf{Classification:} Sulfolane class satisfies the
coordination geometry requirement and the concentration window
requirement.  It partially fails the reductive stability
requirement (sulfolane is marginally stable at Li metal
potentials but SEI chemistry is not well characterised).
Status: \textbf{Reachable in principle; viscosity and melting
point are the practical gating factors, not the framework
mechanism.}

% ─── 4. Missing Class 3: Fluorinated Ether Sub-Classes ───────────────────────
\section{Missing Class 3: Fluorinated Ether Sub-Classes}
\label{sec:fluorinated_ether}

\subsection{HFTHP and BTFMD: The MU Session Findings}

The MU session Query~6 investigated two fluorinated ether
structures specifically designed for LT Li metal batteries:

\begin{itemize}[leftmargin=2em]
  \item \textbf{HFTHP} (1,1,1,3,3,3-hexafluoro-2-(trifluoromethyl)-
        propan-2-ol, or hexafluorinated THF analogue): high
        fluorination, very low donor number, very weak Li$^+$
        coordination.
  \item \textbf{BTFMD} (bis(2,2,2-trifluoroethyl)methyl
        orthoformate-derived ether): partially fluorinated,
        intermediate coordination strength.
\end{itemize}

MU Query~6 result:
\begin{itemize}[leftmargin=2em]
  \item HFTHP at LiFSI 1.0~M: very high CIP/AGG fraction
        (weak solvent coordination allows extensive anion
        participation).  Estimated SCE $\approx 1.65$--$1.75$.
        Above $\mathrm{SCE}^{*}$, approaching the practical ceiling.
  \item BTFMD at LiFSI 1.0~M: moderately high CIP/AGG fraction.
        Estimated SCE $\approx 1.55$--$1.65$.  Significantly above
        $\mathrm{SCE}^{*}$.
\end{itemize}

Both are in the arctic design space (SCE above $\mathrm{SCE}^{*}$).
Both have negative temperature coefficients of SCE (fluorinated
ethers have weaker coordination that strengthens relative to
anion binding as temperature drops, reducing CIP/AGG fraction).

\subsection{HFTHP: Arctic Candidate B (Pre-Registered Prediction)}

HFTHP alone has $\mathrm{SCE} \approx 1.70$, which is too far
above $\mathrm{SCE}^{*}$ for the temperature-responsive crossing
to occur at $-30\,^\circ$C with a realistic temperature
coefficient.  However, HFTHP/DME mixtures can bring the RT SCE
into the arctic design space:

\textbf{Pre-registered prediction (HFTHP/DME arctic candidate B):}
LiFSI 1.0~M in HFTHP/DME at 1:5 v/v:
\begin{itemize}[leftmargin=2em]
  \item Estimated RT SCE $\approx 1.55$--$1.60$.
  \item Estimated temperature coefficient $\approx 0.0018$--$0.0025$
        SCE~units~per~$^\circ$C.
  \item Predicted crossing temperature $T^{*} \approx -20$ to
        $-40\,^\circ$C depending on actual temperature coefficient.
  \item Falsification condition: CE$_\mathrm{LT}$ at $-30\,^\circ$C
        exceeds CE$_\mathrm{LT}$ at $-20\,^\circ$C (non-monotonic
        temperature profile, as in Arctic Candidate~A).
\end{itemize}

\textbf{Classification:} HFTHP satisfies the coordination geometry
requirement (fluorinated ether oxygen coordination is geometrically
distinct from unfluorinated ether coordination) and the
concentration window requirement.  It has been demonstrated
to be stable at Li metal potentials in published literature
\cite{Holoubek2021}.  Status: \textbf{Reachable.  Arctic
Candidate~B pre-registered.}

\subsection{BTFMD: Intermediate Fluorinated Ether}

BTFMD's SCE of $\approx 1.55$--$1.65$ places it in the arctic
design space directly, without requiring a co-solvent.  The
temperature coefficient of BTFMD is expected to be larger than
HFTHP's because BTFMD's partial fluorination produces a
stronger temperature sensitivity in the CIP/AGG equilibrium.

\textbf{Pre-registered prediction (BTFMD pure):}
LiFSI 1.0~M in pure BTFMD:
\begin{itemize}[leftmargin=2em]
  \item Estimated RT SCE $\approx 1.58$--$1.65$.
  \item Predicted crossing temperature $T^{*} \approx -15$ to
        $-35\,^\circ$C.
  \item Falsification condition: same non-monotonic temperature
        profile requirement as Arctic Candidate~A.
\end{itemize}

If confirmed, BTFMD would be the first \emph{single-component}
electrolyte (excluding salt) to exhibit the temperature-responsive
SCE gap without any co-solvent requirement.

\textbf{Classification:} BTFMD satisfies all three requirements.
Status: \textbf{Reachable.  Pre-registered as Arctic Candidate~B
variant.}

% ─── 5. Missing Class 4: Carbonates in Ether Mixtures ────────────────────────
\section{Missing Class 4: Carbonate-Class Solvents in Ether Mixtures}
\label{sec:carbonate}

\subsection{Why Carbonates Are Absent}

Carbonate solvents (EC, DMC, EMC, DEC, FEC) are the dominant
class in Li-\emph{ion} battery electrolytes but are largely
absent from Li \emph{metal} electrolyte optimisation because:
\begin{itemize}[leftmargin=2em]
  \item Carbonates form an inorganic-organic SEI on Li metal
        that suppresses Li plating morphology differently from
        ether-based SEIs.
  \item The carbonate-based SEI is less flexible and more prone
        to cracking under Li metal volume change during cycling.
  \item Carbonates have lower electrochemical stability windows
        than ethers at high voltages but are acceptable at
        moderate voltages.
\end{itemize}

For Li metal anodes, ether-based electrolytes are preferred.
The SCE Framework dataset reflects this preference.

\subsection{SCE Position of Carbonate-Class Solvents}

Carbonate coordination to Li$^+$ is via the carbonyl oxygen
(C=O).  The carbonyl ESP minimum ($\approx -0.75$~eV for DMC)
is more negative than the ether ESP minimum ($\approx -0.62$~eV
for DME) but less negative than the phosphoryl ESP minimum
($\approx -0.79$~eV for TMP).

This places carbonates in the ESP range between ethers and
phosphate esters.  The carbonate C=O--Li$^+$ geometry is
monodentate for DMC and bidentate for EC (the ring carbonate
can chelate through both oxygens in some configurations).

\begin{proposition}[Carbonate ESP position]
Carbonates occupy the ESP range between ethers and phosphate
esters:
\begin{equation}
  |\mathrm{ESP}_{\min}(\mathrm{DME})| <
  |\mathrm{ESP}_{\min}(\mathrm{DMC})| <
  |\mathrm{ESP}_{\min}(\mathrm{TMP})|
\end{equation}
A carbonate/ether mixture can in principle produce cross-class
ESP matching analogous to the DME/TMP mechanism (Preprint~2),
with the carbonate playing the role of TMP.  The ESP gap
between DME and DMC is larger than between DME and TMP
($\approx 0.13$~eV vs.\ $0.007$~eV), reducing the population
splitting effectiveness but not precluding it at higher
carbonate fractions.
\end{proposition}

\subsection{The SEI Chemistry Complication}

The primary obstacle for carbonate/ether mixtures in Li metal
batteries is not the SCE mechanism but the SEI chemistry:
carbonate decomposition at Li metal produces an SEI that
competes with ether-derived SEI chemistry.  The framework
equations (Regime~1) are derived from ether-dominated SEI
systems.  A carbonate/ether mixture at significant carbonate
fractions may produce a hybrid SEI whose CE is not described
by the Regime~1 equations.

\textbf{Pre-registered prediction (DMC/DME, framework scope test):}
LiFSI 1.0~M in DMC/DME at 1:4 v/v:
\begin{itemize}[leftmargin=2em]
  \item Raman SCE estimated at $\approx 1.35$--$1.45$.
  \item If the framework equations apply (ether-dominated SEI):
        CE$_\mathrm{RT}$ predicted $\approx 72$--$74\%$.
  \item If SEI chemistry is dominated by carbonate decomposition:
        CE$_\mathrm{RT}$ may deviate significantly from the
        prediction without falsifying the SCE mechanism.
  \item Diagnostic: FEC-free vs.\ FEC-containing comparison
        to isolate SEI chemistry contribution.
\end{itemize}

\textbf{Classification:} Carbonate class satisfies the
coordination geometry requirement and the concentration window
requirement.  The reductive stability requirement is partially
satisfied (EC is problematic at Li metal; DMC is marginally
stable; FEC forms a stable SEI but introduces a dedicated
SEI-forming mechanism).  The framework equations may require
modification for carbonate-containing systems.  Status:
\textbf{Conditionally reachable.  Framework applicability to
carbonate SEI systems requires dedicated characterisation.}

% ─── 6. Missing Class 5: Ionic Liquid Co-Solvents ────────────────────────────
\section{Missing Class 5: Ionic Liquid Co-Solvents}
\label{sec:ionic_liquid}

\subsection{Why Ionic Liquids Are Absent}

Ionic liquids (ILs) are absent from the confirmed dataset because:
\begin{itemize}[leftmargin=2em]
  \item Most ILs have high viscosity at RT and become gels or
        solids at LT, preventing the low-temperature performance
        measurement that is central to the framework.
  \item IL cations (pyrrolidinium, imidazolium, piperidinium)
        introduce a competing cation that modifies the Li$^+$
        coordination environment in ways not captured by the
        two-component SCE analysis.
  \item The definition of SCE requires a well-defined Li$^+$
        first-shell with a finite number of distinguishable
        coordination configurations.  In IL-rich systems,
        the cation--anion--solvent interactions are highly
        entangled, and the Li$^+$ first shell is not cleanly
        separable from the bulk ionic environment.
\end{itemize}

\subsection{Framework Applicability to IL-Containing Systems}

The SCE Framework in its current form applies to systems where:
\begin{enumerate}[leftmargin=2em]
  \item The Li$^+$ first solvation shell can be defined and
        measured by Raman spectroscopy.
  \item The coordination configurations are classifiable into
        a finite number of distinguishable types.
  \item The SCE of the Li$^+$ shell is separable from the
        bulk electrolyte structural properties.
\end{enumerate}

For IL co-solvent systems at low IL fractions ($< 10\%$ v/v),
these conditions may be approximately satisfied: the IL cation
is present in the bulk but may not enter the Li$^+$ primary
shell at low concentration.  At higher IL fractions, the
cation's influence on Li$^+$ coordination becomes significant
and the framework's two-component model requires extension.

\begin{proposition}[Ionic liquid scope limit]
The SCE Framework in its current form applies to IL/ether
systems only at IL mole fractions where the IL cation does
not measurably enter the Li$^+$ primary coordination shell.
At higher IL fractions, the framework requires extension to
a three-species coordination model (solvent, anion, cation).
This extension is not derived in the present series.
Classification: \textbf{Scope Limit A applies at high IL fraction.
At low IL fraction ($< 5\%$ mol), the framework may apply as
an approximation.}
\end{proposition}

\textbf{Pre-registered prediction (low-IL limit test):}
LiFSI 1.0~M in Pyr$_{14}$FSI (3\% mol)/DME (97\% mol):
\begin{itemize}[leftmargin=2em]
  \item Raman SCE estimated at $\approx 1.26$--$1.30$
        (slightly above pure DME, small CIP contribution
        from additional FSI$^-$ from IL).
  \item CE predictions follow Regime~1 framework equations.
  \item Falsification condition: if Raman shows pyrrolidinium
        ring peaks in the Li$^+$--O stretch region
        ($300$--$500$~cm$^{-1}$), the IL cation is entering
        the primary shell and the framework prediction is
        invalidated for that system.
\end{itemize}

% ─── 7. Missing Class 6: Solid-State and Quasi-Solid Analogues ───────────────
\section{Missing Class 6: Solid-State and Quasi-Solid Electrolyte Analogues}
\label{sec:solid_state}

\subsection{Framework Applicability}

The SCE Framework is derived from liquid-phase electrolyte
systems where the Li$^+$ solvation shell is a dynamic,
thermally accessible distribution of coordination configurations.
Solid electrolytes (oxide, sulfide, halide garnets, polymer)
present a qualitatively different coordination environment:

\begin{itemize}[leftmargin=2em]
  \item Li$^+$ occupies fixed crystallographic sites.
  \item The ``coordination'' is not a dynamic equilibrium
        among distinguishable configurations but a static
        site occupation.
  \item Shannon entropy of coordination configurations is
        replaced by site occupation entropy, which is a
        different quantity with different physical significance.
\end{itemize}

\begin{proposition}[Solid electrolyte scope limit]
The SCE Framework as derived does not apply to solid
electrolytes.  The fixed-point $\mathrm{SCE}^{*} = 1.466$
is defined for liquid-phase dynamic coordination equilibria.
The concept of coordination entropy is meaningful in this
context but the specific numerical value and its relation
to CE cannot be transferred to solid systems without
re-derivation.  Classification: \textbf{Scope Limit A
applies (no dynamic coordination geometry).  Out of scope
in current form.}
\end{proposition}

\subsection{Quasi-Solid and Gel Electrolytes}

Quasi-solid and gel electrolytes (polymer matrices swollen
with liquid electrolyte, PVDF-HFP gels, PEO with plasticisers)
present an intermediate case.  If the liquid fraction retains
a dynamic Li$^+$ solvation equilibrium, the SCE of the
liquid fraction may be definable and measurable by Raman.
In this case, the framework may apply to the liquid fraction
with the understanding that the polymer matrix provides a
modified activity coefficient environment.

\textbf{Pre-registered prediction (PEO-LiFSI/DME quasi-solid):}
A quasi-solid with 30\% PEO (MW 600,000), LiFSI 1.0~M,
DME/TMP 4:1 (liquid fraction):
\begin{itemize}[leftmargin=2em]
  \item If Raman shows a well-resolved Li$^+$--solvent solvation
        structure in the liquid fraction: apply the framework
        to the liquid fraction SCE.
  \item Predicted SCE (liquid fraction only): $\approx 1.42$--$1.50$
        (slightly modified from pure liquid due to PEO ether
        oxygen competition for Li$^+$).
  \item CE prediction follows framework equations for the
        liquid fraction composition.
  \item Falsification: if Raman shows PEO ether oxygens
        dominating the Li$^+$ coordination shell, the system
        requires a modified framework that includes polymer
        coordination explicitly.
\end{itemize}

% ─── 8. The Three Scope Limits ────────────────────────────────────────────────
\section{The Three Scope Limits of the SCE Framework}
\label{sec:scope_limits}

\subsection{Scope Limit A: The Coordination Geometry Requirement}

\begin{definition}[Coordination geometry requirement]
The SCE Framework applies to electrolyte systems where Li$^+$
has access to at least two structurally distinguishable
coordination configurations with non-negligible thermal
population ($p_i \geq 0.05$) under the conditions of interest.
\end{definition}

\textbf{Classes that fail this requirement:}
\begin{itemize}[leftmargin=2em]
  \item Nitrile solvents (geometrically rigid linear
        C$\equiv$N--Li$^+$, no accessible sub-configurations).
  \item Solid electrolytes (static crystallographic sites).
  \item Ionic liquid-rich systems (entangled cation--anion
        coordination obscures the Li$^+$ shell definition).
  \item Any system where Li$^+$ is solvated by a single
        rigid monodentate ligand with no conformational
        or compositional sub-configurations.
\end{itemize}

\subsection{Scope Limit B: The Reductive Stability Requirement}

\begin{definition}[Reductive stability requirement]
The SCE Framework applies to electrolyte systems where the
solvent and salt remain electrochemically stable at Li metal
potentials ($0$--$0.5$~V vs.\ Li/Li$^+$) for at least
100 cycles, allowing the solvation structure equilibrium
to be established and the CE to reflect the geometry-driven
mechanism rather than irreversible solvent decomposition.
\end{definition}

\textbf{Classes that fail this requirement:}
\begin{itemize}[leftmargin=2em]
  \item Nitrile solvents at Li metal potentials (reductive
        polymerisation, irreversible SEI growth).
  \item Carbonate-class solvents at high carbonate fractions
        in Li metal cells (SEI chemistry competes with
        geometry-driven mechanism).
  \item High-voltage additives and co-solvents that are
        reductively unstable below 2~V vs.\ Li/Li$^+$.
\end{itemize}

\subsection{Scope Limit C: The Concentration Window Requirement}

\begin{definition}[Concentration window requirement]
The SCE Framework applies to electrolyte systems where the
accessible concentration range (0.5--2.5~M LiFSI) contains
at least one concentration at which $\mathrm{SCE}
\approx \mathrm{SCE}^{*}$, or where a co-solvent mechanism
can produce $\mathrm{SCE} \approx \mathrm{SCE}^{*}$ at
a chemically stable composition.
\end{definition}

\textbf{Classes that approach this limit:}
\begin{itemize}[leftmargin=2em]
  \item DPE (concentration-invariant, SCE $\approx 1.28$ at
        all accessible LiFSI concentrations, Derivation~6).
        DPE cannot reach $\mathrm{SCE}^{*}$ by concentration
        tuning alone.  It requires a co-solvent.
  \item HFTHP pure (SCE $\approx 1.65$--$1.75$, approaching
        the practical ceiling; concentration tuning reduces SCE
        but may not bring it to $\mathrm{SCE}^{*}$ before
        crossing the practical ceiling from above).
\end{itemize}

% ─── 9. The Regime Boundary and the Band ─────────────────────────────────────
\section{The Regime Boundary and the SCE Band}
\label{sec:band}

\subsection{The Band: Key New Result from the MU Session}

Master Document~8 Part~V identifies the \textbf{Band} as the
key new result of the MU session.  The Band is the SCE range:

\begin{equation}
  1.44 \leq \mathrm{SCE} \leq 1.50
  \label{eq:band}
\end{equation}

Systems with SCE in this range simultaneously satisfy:
CE$_\mathrm{RT} \geq 72\%$ and CE$_\mathrm{LT} \geq 59\%$
from the Regime~1 framework equations.  The Band is the
\emph{target zone} for any formulation aiming for dual
high performance.

The significance of the Band is that it relaxes the precision
requirement on the SCE target: rather than requiring exact
placement at $\mathrm{SCE}^{*} = 1.466$, the framework
specifies a Band within which both CE criteria are
simultaneously met.  The Band width of $0.06$ SCE units
corresponds to a population measurement precision of
$\pm 0.02$ in each configuration fraction --- achievable
by Raman deconvolution in practice.

\subsection{Regime Boundary Position}

From Derivation~7, the regime boundary crossing occurs when:

\begin{equation}
  \mathrm{CE}_\mathrm{RT} > 90\% \quad \Leftrightarrow \quad
  \mathrm{SCE} < 0.71
  \label{eq:regime_boundary}
\end{equation}

from the Regime~1 equation.  However, no system in the dataset
achieves CE$_\mathrm{RT} > 90\%$ in the ether class at
moderate LiFSI concentration.  The regime boundary is relevant
only for concentrated electrolyte systems (4~M LiFSI and above)
where AGG configurations dominate and the Regime~2 mechanism
(concentration-driven) operates.

The SCE Framework's operating domain is Regime~1.  Regime~2 is
a separate design problem not addressed in this series.

% ─── 10. The Complete Novelty Stack ──────────────────────────────────────────
\section{The Complete Novelty Stack of the Series}
\label{sec:novelty_stack}

The following is the complete enumeration of novel contributions
across all six preprints, derived and pre-registered before
experimental confirmation.

\subsection{Novel Theoretical Results}

\begin{enumerate}[leftmargin=2em]
  \item \textbf{The fixed point $\mathrm{SCE}^{*} = 1.466$}
        (Preprint~1): the first analytically derived fixed
        point of Li metal battery electrolyte performance,
        derived from the intersection of two opposing empirical
        gradients.  Not previously identified in any published
        work.
  \item \textbf{The two-gradient structure} (Preprint~1):
        the identification that RT and LT CE respond to SCE
        with opposing gradients, creating a unique performance
        maximum at their intersection.
  \item \textbf{The ESP-matching mechanism} (Preprint~2):
        the identification that co-solvents with ESP minima
        within $k_BT$ of the base solvent produce maximal
        population splitting and maximum SCE increase per
        unit fraction.
  \item \textbf{The anion receptor mechanism} (Preprint~3):
        the identification that selective Lewis acid addition
        can contract the CIP/AGG population and navigate SCE
        from the anion-participation-driven direction.
  \item \textbf{The within-class geometric diversity mechanism}
        (Preprint~4): the identification that cyclic/linear
        ether mixing produces a weaker but practically more
        accessible SCE tuning mechanism within a single chemical
        class.
  \item \textbf{The temperature-responsive SCE gap} (Preprint~5):
        the identification that formulations above $\mathrm{SCE}^{*}$
        at RT with negative temperature coefficients of SCE
        can cross the fixed point at a target operating
        temperature.  The arctic optimum as a design principle.
  \item \textbf{The SCE ceiling} (Preprint~5): the derivation of
        the maximum achievable SCE in ether-based Li metal
        electrolyte systems ($\approx 1.75$) and the arctic
        design space between the ceiling and $\mathrm{SCE}^{*}$.
  \item \textbf{The regime boundary failure mode} (Derivation~7):
        the identification that Arctic designs can enter a
        regime where the Regime~1 equations fail if RT SCE
        exceeds the practical ceiling.
  \item \textbf{The geometric closure of the ether design space}
        (MU Molecular Search): the systematic assignment of every
        ether-class molecule to a region (reached, unreachable,
        discarded) relative to $\mathrm{SCE}^{*}$.
  \item \textbf{The three scope limits} (this paper): the
        formal derivation of Scope Limits A, B, and C defining
        the framework's operating domain.
  \item \textbf{The Band} (MU session, this paper): the
        identification of the target zone $1.44 \leq \mathrm{SCE}
        \leq 1.50$ in which both CE criteria are simultaneously met.
  \item \textbf{The DPE concentration invariance prediction}
        (Derivation~6): the first explicit prediction that DPE
        is concentration-invariant with respect to SCE, providing
        a negative control for the concentration tuning mechanism.
\end{enumerate}

\subsection{Novel Candidate Formulations (Pre-Registered)}

\begin{enumerate}[leftmargin=2em]
  \item Candidate~4: LiFSI 1.2~M DME/TMP 4:1 v/v (Preprint~2).
  \item Candidate~3: LiFSI 1.0~M DOL/TMB 4:1 v/v (Preprint~3).
  \item Candidate~1: LiFSI 1.2~M 2-MeTHF/DME 1.7:1 v/v (Preprint~4).
  \item Candidate~2: LiFSI 1.0~M FEME/2-MeTHF $\approx$0.9:1 v/v
        (Preprint~4).
  \item Arctic Candidate~A: LiFSI 1.2~M DME/TMP/TMB 3:1:0.5
        v/v/v (Preprint~5).
  \item Arctic Candidate~B: LiFSI 1.0~M HFTHP/DME 1:5 v/v
        (this paper, pre-registered).
  \item Arctic Candidate~B variant: LiFSI 1.0~M BTFMD pure
        (this paper, pre-registered).
  \item Sulfolane/DME test: LiFSI 1.0~M sulfolane/DME 1:4 v/v
        (this paper, pre-registered).
  \item Low-IL limit test: LiFSI 1.0~M Pyr$_{14}$FSI (3\% mol)/DME
        (this paper, pre-registered).
  \item Carbonate scope test: LiFSI 1.0~M DMC/DME 1:4 v/v
        (this paper, pre-registered).
\end{enumerate}

\subsection{Novel Experimental Predictions (All Pre-Registered)}

\begin{enumerate}[leftmargin=2em]
  \item Non-monotonic temperature-CE profile for arctic
        candidates (maximum LT CE at $T^{*}$, unique signature
        of arctic optimum).
  \item DPE concentration invariance: $\Delta\mathrm{SCE}/\Delta c
        \approx 0.02$~L~mol$^{-1}$ vs.\ FEME's $0.19$~L~mol$^{-1}$.
  \item FEME crossing concentration: $c^{*} = 1.51$~M for pure
        FEME to reach $\mathrm{SCE}^{*}$.
  \item TEB gradient test: TEB requires higher volume fraction
        than TMB to reach $\mathrm{SCE}^{*}$, confirming Lewis
        acidity dependence of the anion receptor mechanism.
  \item Nitrile preferential coordination: AN coordinates Li$^+$
        preferentially at all fractions above 5\% in DME,
        precluding population splitting.
  \item Sulfolane non-monotonic SCE: sulfolane/DME SCE
        decreases as DME fraction increases, approaching
        $\mathrm{SCE}^{*}$ from above.
  \item BTFMD single-component arctic: BTFMD alone exhibits
        a non-monotonic temperature-CE profile without
        co-solvent addition.
\end{enumerate}

% ─── 11. The Programme for What Remains ──────────────────────────────────────
\section{The Programme for What Remains}
\label{sec:programme}

\subsection{Experimental Priorities}

\paragraph{Priority 1 (Confirmatory, no new candidates):}
Test Candidates~1--4 and Arctic Candidate~A in order of
practical accessibility: Candidate~1 (2-MeTHF/DME, all
reagents standard), Candidate~4 (DME/TMP, Preprint~2),
Candidate~3 (DOL/TMB, Preprint~3), Arctic Candidate~A
(DME/TMP/TMB), Candidate~2 (FEME/2-MeTHF).  Each test
confirms or falsifies specific pre-registered conditions.

\paragraph{Priority 2 (New class confirmation):}
Test sulfolane/DME and HFTHP/DME to confirm or falsify
the missing class predictions.  These tests extend the
dataset and confirm whether the framework applies to
non-ether primary solvents.

\paragraph{Priority 3 (Negative controls):}
DPE concentration invariance scan and AN/DME preferential
coordination test.  These tests confirm the scope limits
directly and provide the negative controls needed to
distinguish the mechanism from trivial effects.

\paragraph{Priority 4 (Arctic Candidate B):}
HFTHP/DME and BTFMD pure temperature scans.  These test
whether the temperature-responsive SCE gap is reproducible
across different fluorinated ether systems, confirming the
generality of the arctic design principle.

\subsection{Computational Priorities}

\paragraph{Computational Priority 1:}
MD simulation of the DME/TMP system (Candidate~4) to
compute the actual population fraction distribution and
confirm the population table in Preprint~2.  This closes
the supplementary population fraction gap identified in
the series assessment.

\paragraph{Computational Priority 2:}
ESP calculation files for TMP and TMB at the DFT level
(B3LYP/6-311+G**, polarisable continuum model for DME
solvent environment).  These constitute the quantum chemistry
audit trail referenced in Preprint~2.

\paragraph{Computational Priority 3:}
Temperature-dependent MD simulation of Arctic Candidate~A
at $+25\,^\circ$C, $0\,^\circ$C, $-20\,^\circ$C,
$-30\,^\circ$C to confirm the temperature-SCE trajectory
and validate the individual temperature coefficient
contributions in Preprint~5 Section~4.

\paragraph{Computational Priority 4:}
Extension of the SCE framework to include the Li$^+$
second solvation shell.  The current framework treats
only the first shell.  Including the second shell may
modify the population fractions and the SCE value at
the fixed point.

\subsection{Theoretical Extensions}

\paragraph{Extension 1 (Multi-cation framework):}
Derive the SCE fixed point for Na$^+$, K$^+$, Zn$^{2+}$,
and Mg$^{2+}$ using the same methodology as Preprint~1.
The Luo et al.\ \cite{Luo2025} paper demonstrates the
universality principle applies across metal ions.  A fixed
point may exist for each cation.

\paragraph{Extension 2 (Regime 2 framework):}
The concentrated electrolyte regime (4~M LiFSI and above)
operates by a different mechanism.  The Regime~2 framework
requires a new derivation analogous to Preprint~1 but
applied to AGG-dominated systems.

\paragraph{Extension 3 (Polymer matrix SCE):}
Extend the framework to gel and quasi-solid electrolytes
by modelling the polymer matrix as an additional coordination
species with a temperature-dependent coordination fraction.

\paragraph{Extension 4 (DSL formulation):}
The OrganismCore DSL (META\_DSL.md, same repository) provides
primitives for expressing the SCE navigation problem as a
formal reasoning object.  A future extension of this series
would encode the full design space navigation --- the fixed
point, the mechanisms, the candidates, the falsification
conditions --- as a machine-executable reasoning substrate.

% ─── 12. Discussion: The Series as a Complete Programme ──────────────────────
\section{Discussion: The Series as a Complete Scientific Programme}
\label{sec:discussion}

\subsection{What Has Been Achieved}

Across six preprints, the SCE Framework has accomplished:

\begin{enumerate}[leftmargin=2em]
  \item \textbf{Derived a fixed point from data.} The fixed point
        $\mathrm{SCE}^{*} = 1.466$ was derived analytically from
        the intersection of two empirical gradients.  It was not
        fitted post-hoc.  It was derived before any candidate
        formulations were specified.
  \item \textbf{Derived three navigation mechanisms.} Each
        mechanism was derived independently from the fixed point,
        not from empirical screening.  Each has specific pre-
        registered predictions.
  \item \textbf{Derived one temperature-responsive design principle.}
        The arctic optimum was derived from the temperature
        dependence of the fixed-point navigation, not from
        empirical observation.
  \item \textbf{Specified ten candidate formulations.} All ten
        are pre-registered with specific falsification conditions
        before any experimental test.
  \item \textbf{Confirmed against six published papers.} None
        falsifies the framework.  Two provide independent
        confirmation.
  \item \textbf{Achieved geometric closure of the ether design space.}
        Every ether-class molecule has been assigned a status
        relative to $\mathrm{SCE}^{*}$.
  \item \textbf{Defined three scope limits.} The framework's
        operating domain is precisely specified.
  \item \textbf{Identified twelve novel theoretical results.}
        None was previously identified in the published literature.
\end{enumerate}

\subsection{The Nature of What Has Been Built}

The series constitutes a \emph{rational design framework} in the
strict sense: a system in which the design target is derived
before experiment, the mechanisms for reaching it are derived
before experiment, the candidate formulations are derived before
experiment, and the falsification conditions are specified before
experiment.

This is not standard practice in electrolyte science.  Standard
practice is:
\begin{enumerate}[leftmargin=2em]
  \item Screen many formulations.
  \item Retain the best performer.
  \item Reason backward about why it worked.
  \item Propose a new screen informed by the retrospective reasoning.
\end{enumerate}

The SCE Framework inverts this process at every step.  The inversion
is made possible by the existence of the fixed point: a
mathematically derived target that converts an unbounded search
into a bounded navigation problem.

\subsection{The Honest Assessment of What Remains Unknown}

Three things remain genuinely unknown until experiment:

\begin{enumerate}[leftmargin=2em]
  \item \textbf{Whether the population fraction tables are correct.}
        The framework predicts specific population fractions at
        each candidate formulation.  These are derived from
        first-principles arguments and MU session estimates.
        They have not been confirmed by experiment or simulation.
        If the tables are systematically wrong by more than
        $\pm 0.05$ in each fraction, the SCE predictions will
        be off by more than the Band width of $0.06$ and the
        ratio predictions will require correction.

  \item \textbf{Whether the 2-MeTHF film-forming contribution
        adds to or modifies the geometry-driven mechanism.}
        The 2-MeTHF dataset point ($\mathrm{SCE} = 1.55$,
        CE$_\mathrm{RT} = 86\%$) exceeds the Regime~1
        geometry-driven prediction.  The excess is attributed
        to ring-opening film-forming chemistry.  Whether this
        contribution persists, disappears, or changes sign
        in the 2-MeTHF/DME mixed system requires experiment.

  \item \textbf{Whether the temperature coefficient estimates
        for Arctic Candidate~A are correct.}  The three-
        contribution additive model and the cooperative
        cooling correction are first-principles estimates
        with $\pm 30\%$ uncertainty.  The crossing temperature
        prediction of $-27\,^\circ$C may be off by as much as
        $\pm 10\,^\circ$C.  The falsification condition P1
        (non-monotonic profile) is robust to this uncertainty;
        the exact crossing temperature is not.
\end{enumerate}

These uncertainties are precisely stated and each has a specific
experimental test.  That is the correct scientific posture.

\subsection{The Repository as the Geometric Backbone}

The decision to make the repository the primary record was not
a formatting choice.  It is a scientific choice: the repository's
timestamped commit history constitutes the pre-registration record.
It is the proof that the framework was derived before experiment.
Without the commit history, the claims of derivation priority are
assertions.  With it, they are verifiable facts with timestamps.

Every derivation, every population fraction table, every
falsification condition, every candidate formulation exists as
a timestamped commit before any experimental test.  The repository
is the audit trail.  The papers are nodes in the graph it describes.

% ─── 13. Conclusion ───────────���──────────────────────────────────────────────
\section{Conclusion}

This paper closes the SCE Framework series geometrically.

The ether-class design space is closed: the molecular search
has assigned every ether-class molecule to a status relative
to $\mathrm{SCE}^{*}$.  The non-ether classes are classified
by their relationship to the three scope limits.  The missing
classes that can reach $\mathrm{SCE}^{*}$ (sulfolane/ether,
fluorinated ether sub-classes) have pre-registered predictions.
The classes that cannot (nitriles, solid electrolytes, IL-rich
systems) have formal scope limit assignments.

The complete novelty stack comprises 12 novel theoretical
results, 10 pre-registered candidate formulations, and 7 novel
experimental predictions.  None was present in the published
literature before the series.

The programme for what remains is explicitly defined: 4
experimental priorities, 4 computational priorities, and 4
theoretical extensions.  Each is bounded, specific, and
connected to the pre-registered predictions by the falsification
conditions.

The fixed point $\mathrm{SCE}^{*} = 1.466$ is the geometric
centre of the programme.  Everything derived in this series
is a consequence of its existence.  Everything in the programme
is a test of whether that consequence holds.

\bigskip
\noindent\textbf{Data and derivation record availability.}
All pre-registered predictions, falsification conditions,
population estimates, scope limit derivations, molecular search
records, and MU session records referenced in this paper are
archived with full timestamped commit history at:\\
\url{https://github.com/Eric-Robert-Lawson/attractor-oncology}\\
The repository commit history constitutes the priority and
pre-registration record for all claims in this paper and
the full series.

\noindent\textbf{Relationship to other papers in this series.}
This is Preprint~6 of the SCE Framework series and the series
capstone.  Preprints~1--5 are available at the same repository.
The full geometric structure connecting all papers, mechanisms,
candidates, scope limits, and the programme is navigable from
the repository index.

\noindent\textbf{Competing interests.}
The author declares no competing financial interests.

\noindent\textbf{Author contributions.}
E.R.L.: conceptualisation, class analysis, scope limit
derivation, novelty stack enumeration, programme definition,
writing.

% ─── Appendix ────────────────────────────────────────────────────────────────
\appendix

\section{Complete Class Status Table}
\label{app:class_table}

\begin{longtable}{p{3.5cm}p{2cm}p{2.5cm}p{6cm}}
\caption{Complete status of all chemical classes relative to
the SCE Framework.  Scope limits A, B, C as defined in
Section~\ref{sec:scope_limits}.}
\label{tab:class_status}\\
\toprule
Class & SCE position & Status &
  Notes \\
\midrule
\endfirsthead
\multicolumn{4}{l}{\small (continued)}\\
\toprule
Class & SCE position & Status & Notes \\
\midrule
\endhead
\bottomrule
\endfoot
\bottomrule
\endlastfoot
Ether (linear, DME) &
  $1.24$ (below SCE*) &
  In design space &
  Preprint~1 dataset baseline \\[3pt]
Ether (cyclic, 2-MeTHF) &
  $1.55$ (above SCE*) &
  In design space &
  Dataset; film-forming contribution \\[3pt]
Ether (cyclic, DOL) &
  $1.41$ (near SCE*) &
  In design space &
  Preprint~3 \\[3pt]
Ether (linear, DPE) &
  $1.28$ (below SCE*) &
  Scope Limit C &
  Concentration-invariant; needs co-solvent \\[3pt]
Ether (fluorinated, FEME) &
  $1.37$ at 1M &
  In design space &
  Preprint~4; concentration-tunable \\[3pt]
Ether (fluorinated, HFTHP) &
  $1.65$--$1.75$ (above SCE*) &
  Arctic design space &
  Arctic Candidate~B \\[3pt]
Ether (fluorinated, BTFMD) &
  $1.55$--$1.65$ (above SCE*) &
  Arctic design space &
  Pre-registered (this paper) \\[3pt]
Phosphate ester (TMP) &
  --- (co-solvent only) &
  Mechanism role &
  ESP-matching co-solvent, Preprint~2 \\[3pt]
Lewis acid (TMB) &
  --- (receptor only) &
  Mechanism role &
  Anion receptor, Preprint~3 \\[3pt]
Nitrile (AN, PN) &
  $\leq 0.693$ &
  Scope Limit A + B &
  Geometric rigidity; reductive instability \\[3pt]
Sulfone (sulfolane) &
  $> \mathrm{SCE}^{*}$ (above) &
  Conditionally reachable &
  Viscosity/mp gating; Scope Limit B partial \\[3pt]
Carbonate (DMC, EC) &
  $\approx 1.35$--$1.45$ &
  Conditionally reachable &
  SEI chemistry complication \\[3pt]
Ionic liquid (rich) &
  Undefined &
  Scope Limit A &
  Cation enters Li$^+$ shell at high fraction \\[3pt]
Ionic liquid (low fraction) &
  $\approx$ ether baseline &
  Conditionally reachable &
  Low-IL limit test pre-registered \\[3pt]
Solid electrolyte &
  Not applicable &
  Scope Limit A &
  Static crystallographic sites \\[3pt]
Quasi-solid / gel &
  Liquid fraction only &
  Conditionally reachable &
  Requires liquid-fraction SCE extraction \\[3pt]
\end{longtable}

\section{The Web Structure of the Series}
\label{app:web_structure}

The six preprints form a directed web, not a linear sequence.
The dependency structure is:

\begin{itemize}[leftmargin=2em]
  \item Preprint~1 (foundation): required by all subsequent papers.
  \item Preprint~2 (ESP matching): required by Preprint~5
        (arctic combination).
  \item Preprint~3 (anion reception): required by Preprint~5
        (arctic combination).
  \item Preprint~4 (within-class): independent of Preprints~2 and~3.
  \item Preprint~5 (arctic): requires Preprints~2 and~3.
  \item Preprint~6 (closure): requires all prior preprints.
\end{itemize}

A reader entering the series at any paper can read that paper
independently, then trace back to Preprint~1 for the foundation.
All papers are self-contained.  All papers reference the same
repository for the derivation record.

% ─── References ──────────────────────────────────────────────────────────────
\begin{thebibliography}{99}

\bibitem{Aurbach2000}
Aurbach, D.\ \textit{et al.}
Review of selected electrode--solution interactions which determine
the performance of Li and Li ion batteries.
\textit{J.\ Power Sources} \textbf{89}, 206--218 (2000).

\bibitem{CRC2024}
Haynes, W.\ M.\ (ed.)
\textit{CRC Handbook of Chemistry and Physics}, 105th edn.
CRC Press, Boca Raton, FL (2024).

\bibitem{Chen2022}
Chen, S.\ \textit{et al.}
High-voltage lithium-metal batteries enabled by controlling the
solvation of fluorinated electrolytes.
\textit{Angew.\ Chem. Int.\ Ed.} \textbf{60}, 3402--3410 (2021).

\bibitem{Goodenough2010}
Goodenough, J.\ B.\ \& Kim, Y.
Challenges for rechargeable Li batteries.
\textit{Chem.\ Mater.} \textbf{22}, 587--603 (2010).

\bibitem{Gutmann1976}
Gutmann, V.
\textit{The Donor--Acceptor Approach to Molecular Interactions.}
Plenum Press, New York (1978).

\bibitem{Holoubek2021}
Holoubek, J.\ \textit{et al.}
Tailoring electrolyte solvation for Li metal batteries cycled at
ultra-low temperature.
\textit{Nat.\ Energy} \textbf{6}, 303--313 (2021).

\bibitem{Liu2021}
Liu, B.\ \textit{et al.}
Trimethyl borate as a multifunctional electrolyte additive for
high-performance lithium metal batteries.
\textit{ACS Appl.\ Mater.\ Interfaces} \textbf{13}, 36305--36313 (2021).

\bibitem{Logan2018}
Logan, E.\ R.\ \textit{et al.}
A study of the physical properties of Li-ion battery electrolytes
containing esters.
\textit{J.\ Electrochem.\ Soc.} \textbf{165}, A21--A30 (2018).

\bibitem{Luo2025}
Luo, Y.\ \textit{et al.}
Solvation configurational entropy governs interfacial kinetics in
metal-ion batteries.
\textit{Joule} \textbf{9}, 210--228 (2025).

\bibitem{Qian2015}
Qian, J.\ \textit{et al.}
High rate and stable cycling of lithium metal anode.
\textit{Nat.\ Commun.} \textbf{6}, 6362 (2015).

\bibitem{Reichardt2011}
Reichardt, C.\ \& Welton, T.
\textit{Solvents and Solvent Effects in Organic Chemistry}, 4th edn.
Wiley-VCH, Weinheim (2011).

\bibitem{Tikekar2016}
Tikekar, M.\ D., Choudhury, S., Tu, Z.\ \& Archer, L.\ A.
Design principles for electrolytes and interfaces for stable
lithium-metal batteries.
\textit{Nat.\ Energy} \textbf{1}, 16114 (2016).

\bibitem{Wan2021}
Wan, G.\ \textit{et al.}
Lithium metal anode operating under lean electrolyte conditions:
a review.
\textit{Adv.\ Mater.} \textbf{33}, 2105713 (2021).

\bibitem{Yu2022}
Yu, Z.\ \textit{et al.}
Rational solvent molecule tuning for high-performance lithium metal
battery electrolytes.
\textit{Nat.\ Energy} \textbf{7}, 94--106 (2022).

\bibitem{Wan2020}
Wan, H.\ \textit{et al.}
Bifunctional interphase-enabled lithium metal batteries.
\textit{Chem.\ Mater.} \textbf{32}, 9036--9046 (2020).

\bibitem{Zhang2020}
Zhang, X.-Q., Cheng, X.-B., Chen, X., Yan, C.\ \& Zhang, Q.
Fluoroethylene carbonate additives to render uniform Li deposits in
lithium metal batteries.
\textit{Adv.\ Funct.\ Mater.} \textbf{27}, 1605989 (2017).

\end{thebibliography}

\end{document}
